% !TEX TS-program = lualatex

%%%
% src::
%   url = https://github.com/pfradin/luadraw/discussions/60
%%%

\documentclass[border = 3pt]{standalone}

%%%
% Des couleurs faciles d'emploi via le package \xcolor.
%%%
\usepackage[svgnames]{xcolor}

%%%
% La bibliothèque \luadraw allie une facilité d’utilisation
% à un rendu particulièrement soigné.
%%%
\usepackage[3d]{luadraw}

\begin{document}

\begin{luadraw*}{name = octahedron-in-cube}
local poly = require "luadraw_polyhedrons"

local ld = luadraw

local Origin = ld.pt3d.Origin
local M      = ld.pt3d.M

local vecI = ld.pt3d.vecI
local vecJ = ld.pt3d.vecJ
local vecK = ld.pt3d.vecK

------
-- Usage d'une ¨biblio spécialisée pour récupérer des polyèdres
-- réguliers clés en main (convexes ou non).
------
require 'luadraw_polyhedrons'

------
-- ¨Def de la zone graphique.
------
local graphview = ld.graph3d:new{
  viewdir  = {140, 70},
  adjust2d = true,
  bbox     = false
}

------
-- Les éléments géométriques.
------
local c = 5

local cube = ld.parallelep(
  Origin,
  c*vecI, c*vecJ, c*vecK
)

local shift_vect = -4.5*(vecJ + vecI)

------
-- Dans ''poly.octahedron(C, S)'', les ¨args ''C'' et ''S''
-- désignent le centre et un sommet de l'octaèdre.
------
local octahedron_in_cube = ld.shift3d(
  poly.octahedron(
    Origin,
    (c / 2)*vecK  -- Ceci est un point, et non un vecteur!
  ),
  M(c, c, c) / 2  -- Ceci est un vecteur, et non un point!
)

------
-- Réglages liés aux lignes.
------
graphview:Linewidth(4)

ld.Hiddenlines     = true
ld.Hiddenlinestyle = "dashed"

------
-- Tracé du cube.
------
graphview:Dpoly(
  cube,
  {
-- ''ld.mShadedHidden'' demande la coloration des facettes selon
-- leur inclinaison et l'affichage des arêtes cachées.
    mode  = ld.mShadedHidden,
    color = "LightBlue"
  }
)

------
-- Tracé de deux octaèdres.
------
graphview:Dpoly(
  ld.shift3d(
    octahedron_in_cube,
    shift_vect
  ),
  {
    mode  = ld.mShaded,
    color = "orange"
  }
)

graphview:Dpoly(
  octahedron_in_cube,
  {
-- ''ld.mShadedHidden'' permet d'avoir un rendu de type fil de fer.
    mode      = ld.mWireframe,
    edgestyle = "dashed"
  }
)

------
-- Montrons le résultat de notre oeuvre.
------
graphview:Show()
\end{luadraw*}

\end{document}
